\documentclass[10pt,a4paper]{amsart}
\usepackage[margin=19mm]{geometry}
\usepackage{amsmath,amssymb,mathtools}
\usepackage[scr=rsfs,cal=euler]{mathalfa}
\usepackage{xfrac}
\usepackage[unicode,hidelinks,bookmarks=false]{hyperref}
\newcommand{\ie}{i.e.\ }
\newcommand{\cf}{cf.\ }
\newcommand{\resp}{respectively\ }
\newcommand{\Subsection}{\subsection}
\providecommand{\nobreakdash}{\nobreak\hbox{-}}
\setlength{\emergencystretch}{2em}
\multlinegap0pt
\usepackage{amssymb}
\usepackage{cancel}
\usepackage{xcolor}
\usepackage{siunitx}
\usepackage{tikz}
\usepackage{mathtools}
\newtheorem{proposition}{Proposition}
\newtheorem{lemma}{Lemma}
\newtheorem{theorem}{Theorem}
\def\N{\mathbb{N}}
\def\Z{\mathbb{Z}}
\title{Flexible GMRES converges in two phases}
\author{Stefan G\"uttel, Lauri Nyman}
\date{Source: arXiv:2604.27971v1 (CC BY 4.0)}
\begin{document}
\maketitle

\noindent\emph{Source and licence.} Flexible GMRES converges in two phases. Version arXiv:2604.27971v1 (CC BY 4.0), distributed by arXiv under the Creative Commons Attribution 4.0 International licence, http://creativecommons.org/licenses/by/4.0/, as recorded in the arXiv licence metadata for this exact version and shown on the abstract page https://arxiv.org/abs/2604.27971v1. Full licence notice: \texttt{licenses/arxiv-2604.27971-CC-BY-4.0.txt}.\par\smallskip

\noindent\textit{Editorial scope.} Counted unit: the article's theorem thm:explicit\_bound (source0.tex lines 434--444). The article's complete original proof of it is delivered with it (source0.tex lines 445--503, one \texttt{proof} environment of 59 lines). No mathematics is written, repaired or re-derived: the statement and the proof are the article's own lines, in the article's own notation, and no retained line is substituted or re-flowed.  Retained uncounted as the source-local context the proof needs: lines 360--370; lines 388--410; lines 414--416; lines 423--425.  Nothing was omitted from any retained span, so the package carries no omission record.  No retained line contains a citation, so the wrapper carries no bibliography and no bibliography entry.  No retained line contains a figure environment, an image-inclusion command or drawing code, so no image is delivered and the package declares no asset dependency.  Editorial formatting only, in addition: an AMS \texttt{amsart} preamble that restates the article's own declarations and macros used by the retained lines, namely the environments \texttt{proposition}, \texttt{lemma}, \texttt{theorem} on separate counters, together with the standard packages those lines need.  The retained environments print in this document as Proposition 1, Lemma 1, Theorem 1 in order of appearance, whereas the article prints the same items with the numbering of its own full document.  The complete native source of the article as distributed in the arXiv e-print for this exact version is bundled unchanged as \texttt{sources/199515/source0.tex}.
\begin{proposition} \label{prop:recurrence_gamma}
    It holds that
    \begin{align}
        \begin{aligned} \label{eq:successive_rG_gamma}
             & \left\|r_m^{\mathrm{FG}}\right\| \leq \gamma_m \left\|r_{m-1}^{\mathrm{FG}}\right\|, \\
             & \gamma_m := \frac{\|r_m^{\mathrm{P}}\|}{\sqrt{1-\gamma_{m-1}^2}}, \quad \gamma_0 := 0,
        \end{aligned}
    \end{align}
    for all $m \in \Z^+$ for which $\max_{j \leq m} \gamma_{j}
        < 1$.
\end{proposition}

\subsection{An explicit bound on the FGMRES residual}
The assumption on the~$\gamma_j$ in Proposition~\ref{prop:recurrence_gamma}, namely that $\max_{j \leq m} \gamma_{j} < 1$, guarantees that the denominator in the expression for $\gamma_{j}$ is a positive real number for all $j \leq m$. If we assume 
\[
\left\|r_j^{\mathrm{P}}\right\| = \| v_j - A z_j \| \leq \mu \leq \frac{1}{2} \quad \text{for all $j$},
\]
we can derive a similar recursive bound without needing to assume $\max_{j \leq m} \gamma_{j} < 1$. The resulting bound is easier to analyze and allows us to express a bound for $\left\|r_m^{\mathrm{FG}}\right\|$ purely in terms of the contraction factor~$\mu$.
\begin{lemma} \label{prop:recurrence}
    Let $0 < \mu \leq \frac{1}{2}.$ We have 
    \begin{align}
        \begin{aligned} \label{eq:successive_rG}
             & \left\|r_m^{\mathrm{FG}}\right\| \leq \omega_m \left\|r_{m-1}^{\mathrm{FG}}\right\|, \\
             & \omega_m := \frac{\mu}{\sqrt{1-\omega_{m-1}^2}}, \quad \omega_0 := 0,
        \end{aligned}
    \end{align}
    for all $m \in \Z^+$. 
\end{lemma}
\begin{proof}
    First, we show that $\omega_{j} < 1$ for all $j$ whenever $0 < \mu \leq \frac{1}{2}.$ We do this by showing that if $\omega_j$ lies in the interval $\left(0, \frac{1}{\sqrt{2}}\right]$, so does $\omega_{j+1}$. Indeed, if $0 < \mu \leq \frac{1}{2}$ and $0 < \omega_{j} \leq \frac{1}{\sqrt{2}}$, then
    \begin{align*}
       \omega_{j+1} = \frac{\mu}{\sqrt{1-\omega_{m}^2}} \leq \sqrt{2} \mu \leq \frac{1}{\sqrt{2}}.
    \end{align*}
    This, combined with the fact that clearly $\omega_{j+1} > 0$, proves the result. The remaining proof of the lemma is similar to the proof of Proposition \ref{prop:recurrence_gamma}.
\end{proof}

\begin{lemma} \label{lem:monotonic_omega}
 Let $0 < \mu \leq \frac{1}{2}.$ The sequence $(\omega_j)_{j \in \Z^+}$ is monotonically increasing.
\end{lemma}

\begin{lemma} \label{lem:monotonic_in_mu}
    Let $m \in \Z^+$ and $0<\mu_1 < \mu_2 \leq \frac{1}{2}$. It holds that $\omega_m(\mu_1) < \omega_m(\mu_2)$.
\end{lemma}

\begin{theorem} \label{thm:explicit_bound}
    Let $0 < \mu \leq \frac{1}{2}$. For all $m<N$, the FGMRES residual at step $m$ satisfies
    \begin{align} \label{eq:cheb_bound}
        \left\|r_m^{\mathrm{FG}}\right\| & \leq \left\|r_0\right\| \frac{\mu^{m/2}}{\sqrt{U_m\left(\frac{1}{2\mu}\right)}},
    \end{align}
    where $U_m$ denotes the degree~$m$ Chebyshev polynomial of the second kind.
    Written explicitly, the bound becomes 
    \begin{align} \label{eq:explicit_bound}
        \left\|r_m^{\mathrm{FG}}\right\| & \leq \left\|r_0\right\| \mu^m \left( \frac{\sqrt{1 - 4\mu^2}}{\left( \frac{1 + \sqrt{1 - 4\mu^2}}{2} \right)^{m+1} - \left( \frac{1 - \sqrt{1 - 4\mu^2}}{2} \right)^{m+1}} \right)^{1/2}.
    \end{align}
\end{theorem}

\begin{proof}
    Applying \eqref{eq:successive_rG} recursively leads to
    \begin{align} \label{eq:bound}
        \left\|r_m^{\mathrm{FG}}\right\| \leq \left\|r_0\right\| \prod_{j=1}^m \omega_j.
    \end{align}
    We seek to express $ \prod_{j=1}^m \omega_j$ purely in terms of the contraction factor $\mu$. To simplify the notation, define $x_m := \omega_m^2$ so that the recurrence \eqref{eq:successive_rG} becomes
    $$x_m = \frac{\mu^2}{1 - x_{m-1}}, \quad x_0 = 0.$$
    We approach this fractional recurrence relation by expressing it as a system of two linear recurrence relations. To this end, let $\{a_j\}_{j \in \N}, \{b_j\}_{j \in \N}$ denote sequences such that $x_m = \frac{a_m}{b_m}$. Since $x_0 = \omega_0^2 = 0$, we can set $a_0 = 0$ and $b_0 = 1$. Substituting $x_m = \frac{a_m}{b_m}$ into the recurrence relation gives
    \begin{align} \label{eq:recurrence}
        \frac{a_m}{b_m} = \frac{\mu^2}{1 - \frac{a_{m-1}}{b_{m-1}}} = \frac{\mu^2 b_{m-1}}{b_{m-1} - a_{m-1}}.
    \end{align}
    From this, we define the following system of two linear recurrence relations with initial conditions $a_0 = 0$ and $b_0 = 1$, whose solution clearly satisfies the relation \eqref{eq:recurrence}:
    \begin{align} \label{eq:recurrence_system}
        \begin{aligned}
            a_m & = \mu^2 b_{m-1},     \\
            b_m & = b_{m-1} - a_{m-1}.
        \end{aligned}
    \end{align}
    Now, $x_j = \frac{a_j}{b_j} = \frac{\mu^2 b_{j-1}}{b_j}$, which leads to convenient cancellations in the product $\prod_{j=1}^m x_j$:
    $$\prod_{j=1}^m x_j = \prod_{j=1}^m \left( \frac{\mu^2 b_{j-1}}{b_j} \right) = \mu^{2m} \left( \frac{b_0}{b_1} \cdot \frac{b_1}{b_2} \cdots \frac{b_{m-1}}{b_m} \right) = \mu^{2m} \frac{b_0}{b_m}.$$
    Since $b_0=1$, we have that $\prod_{j=1}^m \omega_j = \frac{\mu^m}{\sqrt{b_m}}.$
    Now, we only need an explicit formula for $b_m$. By substituting the first equation of \eqref{eq:recurrence_system} into the second, we get a single second-order linear recurrence relation for $b_m$:
    \begin{align} \label{eq:2nd_order_recurrence}
        b_m = b_{m-1} - \mu^2 b_{m-2}.
    \end{align}
    The characteristic equation for $b_m - b_{m-1} + \mu^2 b_{m-2} = 0$ is
    $$r^2 - r + \mu^2 = 0,$$
    whose roots are
    \begin{align*}
        r_{\pm} = \frac{1 \pm \sqrt{1 - 4\mu^2}}{2}.
    \end{align*}
    Hence, all solutions are of the form $b_m = C_1(r_+)^m + C_2(r_-)^m$, and using the initial conditions $b_0 = 1$ and $b_1 = 1$, we can solve for the coefficients $C_1, C_2$. Standard algebraic manipulation yields
    \begin{align} \label{eq:bm}
        b_m = \frac{(r_+)^{m+1} - (r_-)^{m+1}}{r_+ - r_-}.
    \end{align}
    This lets us express the bound \eqref{eq:bound} as
    \begin{equation*}
        \begin{aligned}
            \left\|r_m^{\mathrm{FG}}\right\| & \leq \left\|r_0\right\| \mu^m \left(  \frac{r_+ - r_-}{(r_+)^{m+1} - (r_-)^{m+1}} \right)^{1/2}                                                                                       \\
                                             & =\left\|r_0\right\| \mu^m \left( \frac{\sqrt{1 - 4\mu^2}}{\left( \frac{1 + \sqrt{1 - 4\mu^2}}{2} \right)^{m+1} - \left( \frac{1 - \sqrt{1 - 4\mu^2}}{2} \right)^{m+1}} \right)^{1/2}.
        \end{aligned}
    \end{equation*}
    We also note that, after the change of variables $c_m = \frac{b_m}{\mu^m}, \ \mu = \frac{1}{2x}$, the recurrence relation \eqref{eq:2nd_order_recurrence} becomes that of the Chebyshev polynomials of the second kind, which are polynomials defined by 
    $$
    { U_{m}(\cos \theta )\sin \theta ={\sin }{\big (}(m+1)\theta {)}, \quad  0 \leq \theta \leq \pi.}
    $$ As such, we can express the bound \eqref{eq:explicit_bound} alternatively as
    \begin{equation*} 
        \begin{aligned}
            \left\|r_m^{\mathrm{FG}}\right\| & \leq \left\|r_0\right\| \frac{\mu^{m/2}}{\sqrt{U_m\left(\frac{1}{2\mu}\right)}}.
        \end{aligned}
    \end{equation*}
    These bounds are valid at iteration $m$ if the sequence of contraction factors $(\omega_j)_j$ satisfies $\omega_j < 1$ for all $j < m$, and $\omega_m \leq 1$. Setting $\omega_m = 1$ yields $a_m = b_m$, which in turn implies that $b_{m+1}=0.$ This is attained precisely at the roots of $U_{m+1}$, which are well known to be ${x_{k}=\cos \left({\frac {k}{m+2}}\pi \right),\ k=1,\dots ,m+1.}$ Hence, the smallest positive value for $\mu$ for which $\omega_m = 1$ is
    \begin{equation}\label{eq:mum}
        \mu_m = \frac{1}{2 \cos\left(\frac{\pi}{m+2}\right)}.
    \end{equation}
    Lemma~\ref{lem:monotonic_in_mu} implies that, for $\mu < \mu_m$, it holds that $\omega_j < 1$ for all $j \leq m$. %It follows that the range for which the bound \eqref{eq:explicit_bound} holds is $\mu \in [0,\mu_m]$.
    Moreover, it holds that the sequence $(\mu_m)_m$ is monotonically decreasing with $\lim_{m \rightarrow \infty} \mu_m = \frac{1}{2}$. % which coincides with the singularity of \eqref{eq:explicit_bound}. %The monotonicity result of Lemma \ref{lem:monotonic_omega} then implies that if $\mu \in [0,\frac{1}{2}],$ it holds that $\omega_j < 1$ for all $j\in \Z^+$. 
    It follows that the bound \eqref{eq:explicit_bound} holds for $\mu \in [0,\frac{1}{2}].$ This completes the proof.
\end{proof}

\end{document}
